\documentclass[10pt,a4paper]{amsart}
\usepackage[margin=19mm]{geometry}
\usepackage{amsmath,amssymb,mathtools}
\usepackage[scr=rsfs,cal=euler]{mathalfa}
\usepackage{xfrac}
\usepackage[unicode,hidelinks,bookmarks=false]{hyperref}
\newcommand{\ie}{i.e.\ }
\newcommand{\cf}{cf.\ }
\newcommand{\resp}{respectively\ }
\newcommand{\Subsection}{\subsection}
\providecommand{\nobreakdash}{\nobreak\hbox{-}}
\setlength{\emergencystretch}{2em}
\multlinegap0pt
\usepackage{bm}
\usepackage{bbm}
\usepackage{dsfont}
\usepackage{xspace}
\usepackage{cancel}
\usepackage{mathtools}
\usepackage{amsthm}
\makeatletter
\@ifundefined{corollary}{\newtheorem{corollary}{Corollary}}{}
\@ifundefined{theorem}{\newtheorem{theorem}{Theorem}}{}
\makeatother
\title[Granger Causality in Extremes]{Granger Causality in Extremes}
\author{Juraj Bodik, Olivier C. Pasche}
\date{Source: arXiv:2407.09632v3 (CC BY 4.0)}
\begin{document}
\maketitle

\noindent\emph{Source and licence.} Granger Causality in Extremes. Version arXiv:2407.09632v3 (CC BY 4.0), distributed under the Creative Commons Attribution 4.0 International licence, as recorded in the arXiv licence metadata for this exact version and shown on the abstract page https://arxiv.org/abs/2407.09632v3. Full rights notice: licenses/arxiv-2407.09632-CC-BY-4.0.txt.\par\smallskip

\noindent\textit{Editorial scope.} Counted unit: the Corollary whose source label is \texttt{consequence\_1+gamma/2} in the article (source0.tex lines 3466--3488, source label \texttt{consequence\_1+gamma/2}), delivered together with the article's own complete original proof of it (source0.tex lines 3489--3608, one \texttt{proof} environment of 120 lines). No mathematics is written, repaired or re-derived: statement and proof are the article's own text line for line. Required context retained uncounted: sre2 (lines 600--618); Gamma\_hat\_equation (lines 676--678); Theorem3 (lines 751--757). Every definition, display and label of those lines is retained unchanged. Omitted from this package and not redistributed: 1--599, 619--675, 679--750, 758--3465, 3609--4424. Nothing inside a retained excerpt is omitted or altered: every retained body is delivered exactly as the article's lines are written, so no omission receipt of any kind accompanies this package. Citations: the retained excerpts carry no citation command at all, so no bibliography entry of the article is transcribed and no reference is redistributed. Reference adaptation: none was needed and none was made; every reference inside the delivered unit resolves inside it, so no unresolved reference or citation is produced. Printed slips kept literal: no printed word, symbol, hypothesis or number of the counted statement or of its proof was altered. Layout and class: the article's own class is \texttt{article} with packages including geometry, etoolbox, booktabs, xcolor, stackengine, natbib, fontenc, inputenc, bm, bbm, graphicx, subcaption, amsmath, amsfonts; the wrapper is the lane's pinned \texttt{amsart} editorial wrapper at 10pt on A4, which restates the theorem-like environments the retained text needs and adds the article's own macro definitions that the retained text needs (none), with extra wrapper packages (bm, bbm, dsfont, xspace, cancel, mathtools). The delivered numbering is the unit's own. Assets: no retained span contains a figure environment, a graphics-inclusion command, a PDF-page inclusion, a table or a TikZ picture; no image file is copied into this package, the package declares no asset dependency and no image file is redistributed with it. Licence: version arXiv:2407.09632v3 of this article is distributed under the Creative Commons Attribution 4.0 International licence, as recorded in the arXiv licence metadata for this exact version and shown on the abstract page https://arxiv.org/abs/2407.09632v3. A full rights notice is delivered at licenses/arxiv-2407.09632-CC-BY-4.0.txt. Characters: every retained excerpt is pure ASCII, so the delivered engine, XeLaTeX with its default Latin Modern text font, reports no missing character.
\begin{equation}\label{sre2}
\textbf{W}_t = \begin{pmatrix}
Z_t  \\
X_t  \\
Y_t \\
\end{pmatrix}
, 
\textbf{A}_t = \begin{pmatrix}
A^z_{1,t} & A^z_{2,t} & A^z_{3,t} \\
A^x_{1,t} & A^x_{2,t} & A^x_{3,t} \\
A^y_{1,t}& A^y_{2,t} & A^y_{3,t} \\
\end{pmatrix}
, \textbf{B}_t=
\begin{pmatrix}
B^z_{t} \\
B^x_{t}  \\
B^y_{t}  \\
\end{pmatrix}, \varepsilon_t^\cdot = (A^\cdot_{1,t},A^\cdot_{2,t},A^\cdot_{3,t},B^\cdot_t)^\top.
\end{equation}

  \begin{equation}\label{Gamma_hat_equation}
\hat{\Gamma}_{\textbf{X}\to\textbf{Y}\mid\mathcal{\textbf{Z}}}:=\frac{1}{|S|}\sum_{t\in S} F(y_{t+1}),
  \end{equation}

\begin{theorem}\label{Theorem3}
Consider a time series following the non-negative SRE model \eqref{sre2} that satisfies Assumptions (B1), (B2), and (B4). Then, the estimator \(\hat{\Gamma}_{\textbf{X}\to\textbf{Y}\mid \textbf{Z}}\) defined in equation (\ref{Gamma_hat_equation}), with \(S \equiv S_1\) satisfies
\begin{equation}
\label{consistency_theorem2}
\hat{\Gamma}_{\textbf{X}\to\textbf{Y}\mid\textbf{Z}} \overset{P}{\to} 1 \text{   as }n\to\infty\iff  \Gamma_{\textbf{X}\to\textbf{Y}\mid\mathcal{C}} = 1.
\end{equation}
\end{theorem}

\begin{corollary}
\label{consequence_1+gamma/2}
Let the assumptions of Theorem~\ref{Theorem3} hold and suppose that
\(\Gamma_{X\to Y\mid C}<1\). Write 
\[
\widehat\Gamma_n(\tau)
:=
\frac1{|S_1|}\sum_{t\in S_1}F(Y_{t+1}),
\qquad
\widehat\Gamma^{\mathrm{baseline}}_n(\tau)
:=
\frac1{|\widetilde S_1|}
\sum_{t\in \widetilde S_1}F(Y_{t+1}), 
\]
where $\widetilde S_1:=\left\{t\leq n: \begin{pmatrix}Y_t\\ \mathbf Z_t\end{pmatrix}\le \tau\right\}.$ Then there exists \(\tau^0\) such that, for every admissible \(\tau\leq \tau^0\):
\[
P\left(
\widehat\Gamma_n(\tau)
\leq
\frac{1+\widehat\Gamma^{\mathrm{baseline}}_n(\tau)}{2}
\right)\to 1 .
\]
\end{corollary}

\begin{proof}
Recall the notation
\[
D_t(\tau):=\left\{\begin{pmatrix}Y_t\\ \mathbf Z_t\end{pmatrix}\le \tau\right\},
\qquad
S_1:=\{t\le n:X_t\ge \tau_n^X,\ D_t(\tau)\},
\qquad
\widetilde S_1:=\{t\le n:D_t(\tau)\}.
\]
By the characterization used in the proof of Theorem~\ref{Theorem3}, under
(B2) and (B4),
\[
\Gamma_{X\to Y\mid C}<1
\quad\implies \quad
A^y_{2,t+1}=0 \ \text{a.s.}
\]
Therefore, for every  \(t\in D_t(\tau)\), and using (B4), we can write
\[
Y_{t+1}=A^y_{1,t+1}\mathbf Z_t+A^y_{3,t+1}Y_t+B^y_{t+1}\le
A^y_{1,t+1}\tau_Z+A^y_{3,t+1}\tau_Y+B^y_{t+1}.
\]
Since \(F\) is nondecreasing, we always have
\[
\widehat\Gamma_n
\le
U_n(\tau)
:=
\frac1{|S_1|}\sum_{t\in S_1}
K_{t+1}(\tau), \quad \text{ where } \quad K_{t+1}(\tau)
:=
F\!\left(
A^y_{1,t+1}\tau_Z+A^y_{3,t+1}\tau_Y+B^y_{t+1}
\right).
\]
The random variable \(K_{t+1}(\tau)\) is a bounded measurable function of
\(\varepsilon^y_{t+1}\). Thus, by Claim 2 used in the
proof of Theorem~\ref{Theorem3},
\[
U_n(\tau)\overset{P}{\longrightarrow}
e(\tau)
:=
\mathbb E\!\left[
F\!\left(
A^y_{1,t}\tau_Z+A^y_{3,t}\tau_Y+B^y_t
\right)
\right].
\]
It remains to choose \(\tau\) so that \(e(\tau)\) is strictly below the desired midpoint. 

Let \(s_Y:=\inf\operatorname{supp}(Y_t)\) and
\(s_Z:=\inf\operatorname{supp}(\mathbf Z_t)\). In the non-negative SRE case,
\(s_Y\) and \(s_Z\) are finite. Define
\[
e_0
:=
\mathbb E\!\left[
F\!\left(
A^y_{1,t}s_Z+A^y_{3,t}s_Y+B^y_t
\right)
\right].
\]
Since \(F(x)<1\) for every finite \(x\), we have \(e_0<1\). Moreover, by
right-continuity and monotonicity of \(F\), together with dominated convergence,
\[
e(\tau)\longrightarrow e_0
\qquad\text{as}\qquad
\tau\downarrow (s_Y,s_Z).
\]
Hence we may choose \(\tau^0\) sufficiently close to \((s_Y,s_Z)\) such that for every admissible
\(\tau\le \tau^0\),
\[
e(\tau)\le e(\tau^0)<\frac{1+e_0}{2}.
\]
By the ergodic theorem,
\[
\widehat\Gamma^{\mathrm{baseline}}_n(\tau)
=
\frac{\sum_{t=1}^n \mathbf 1\{D_t(\tau)\}F(Y_{t+1})}
     {\sum_{t=1}^n \mathbf 1\{D_t(\tau)\}}
\overset{P}{\longrightarrow}
\gamma_{\mathrm{base}}(\tau):=
\mathbb E\{F(Y_{t+1})\mid D_t(\tau)\}
\]
for every admissible \(\tau\). Again using non-negativity and monotonicity of \(F\),
\[
F(Y_{t+1})
\ge
F\!\left(
A^y_{1,t+1}s_Z+A^y_{3,t+1}s_Y+B^y_{t+1}
\right).
\]
Since \(\varepsilon^y_{t+1}\) is independent of the past, the right-hand side is
independent of \(D_t(\tau)\). Consequently, \(
\gamma_{\mathrm{base}}(\tau)\ge e_0.
\) Thus, for every admissible \(\tau\le\tau^0\),
\[
e(\tau)
<
\frac{1+e_0}{2}
\le
\frac{1+\gamma_{\mathrm{base}}(\tau)}{2}.
\]
Combining this with the convergence and
\(U_n(\tau)\overset{P}{\to}e(\tau)\) and the strict population inequality above gives
\[
P\left(
U_n(\tau)\le
\frac{1+\widehat\Gamma^{\mathrm{baseline}}_n(\tau)}{2}
\right)\to 1.
\]
Since \(\widehat\Gamma_n(\tau)\le U_n(\tau)\), we conclude that also
\(
P\left(
\widehat\Gamma_n(\tau)
\le
\frac{1+\widehat\Gamma^{\mathrm{baseline}}_n(\tau)}{2}
\right)\to 1.
\)
This proves the claim.
\end{proof}

\end{document}
